\section{Fourier-Transformation (+)}\label{MA2-fouriertransformation}

\begin{table}[H]
  \centering
  \begin{tabularx}{\textwidth}{@{}>{\raggedright\arraybackslash}X r@{}}
    \toprule
    \multicolumn{2}{c}{\textbf{�bersicht �ber Lernziele und Voraussetzungen}} \\
    \midrule
    \textbf{Lernziele} & \\[.5\normalbaselineskip]
    \tabitem die Fourier-Transformierte einer absolut integrierbaren Funktion berechnen & \\
    \tabitem grundlegende Eigenschaften der Fourier-Transformation anwenden & \\
    \tabitem die Bedeutung der Fourier-Transformation im Frequenzraum erl�utern & \\[.5\normalbaselineskip]
    \midrule
    \textbf{Voraussetzungen} &  \\[.5\normalbaselineskip]
    \tabitem uneigentliche Integrale berechnen k�nnen & \abschnittsverweis{MA2-integration}  \\
    \tabitem Grenzwerte sicher bilden k�nnen & \abschnittsverweis{MA1-grenzwerte-von-folgen-und-funktionen}  \\[.5\normalbaselineskip]
    \bottomrule
  \end{tabularx}
\end{table}

\noindent \ifintegriereaufgaben Die �bungsaufgaben zum Thema finden Sie \hyperref[MA2-fouriertransformation-aufgaben]{hier}.
\vspace{0.3cm}
\hrule
\fi

\subsection{Fragestellung}

Kann man auch f�r Funktionen, die nicht periodisch, sind ein Spektrum finden? Wenn ja, wie sieht das Spektrum aus und wie berechnet man es? Mit diesen Fragen werden wir uns in diesem Abschnitt besch�ftigen. Die Antwort auf die Frage lautet "\,ja, kann man"\, und das Ergebnis ist eine Funktion, die sogenannte Fourier-Transformierte.

\subsection{Definition}

Betrachten wir eine Funktion \(f: \mathbb R \to \mathbb C, t \mapsto f(t)\), die nicht periodisch ist, aber f�r die gilt:
\begin{itemize}
	\item[(a)] $f$ ist stetig differenzierar und
	\item[(b)] \(f\) ist absolut integrierbar  \(\quad \displaystyle \int\limits_{-\infty}^\infty |f(t)|\,\mathrm{d}t < \infty \,. \)
\end{itemize}

Mit dem folgenden Trick n�hern wir uns langsam der Frage nach dem Spektrum von \(f\):
\begin{enumerate}
	\item Wir w�hlen eine positive, reelle Zahl $T>0$ und fixieren das Intervall $\left(-\frac{T}{2}, \frac{T}{2}\right)$.
	\item Wir definieren eine periodische Funktion $g$, \(g(t) = g(t+T)\),  die \(f\) periodisch fortsetzt:
	\[ \forall t\in \left(-\frac{T}{2}, \frac{T}{2}\right): \; g(t) = f(t)\,.  \]
\end{enumerate}

Die Funktion $g$ erf�llt alle Voraussetzung aus dem vorigen Abschnitt und wir erhalten eine Darstellung von \(g\) als Fourierreihe, insbesondere gilt f�r alle \(t\in \left(-\frac{T}{2}, \frac{T}{2}\right) \):
\[\begin{array}{l rcl}
		-\frac{T}{2} < t < \frac{T}{2}:  & g(t) &=& \displaystyle\sum\limits_{k=-\infty}^\infty \textcolor{blue}{c_k} \cdot  e^{i \, \omega_k \, t} \\
	\textnormal{Formel f�r } c_k: & &=& \displaystyle\sum\limits_{k=-\infty}^\infty \left(\textcolor{blue}{\dfrac1T \displaystyle \int\limits_{-\frac{T}{2}}^{\frac{T}{2}} f(t) \cdot  e^{-i \, \omega_k \, t}\, \mathrm{d}t} \right) \cdot  e^{i \, \omega_k \, t}, \quad \omega_k = k\cdot \dfrac{2\pi}{T}\\
 \textnormal{Relation zu }f: & &=& f(t) \end{array}\]

 Das letzte Gleichheitszeichen gilt per Konstruktion: \(f\) stimmt auf \( \left(-\frac{T}{2}, \frac{T}{2}\right)\) ja mit \(g\) �berein! Naja, wenn das so ist, dann schicken wir \(T\) gegen Unendlich und erhalten eine eine Darstellung von \(f\) f�r alle \(t\in \mathbb R\), oder? Die Idee ist genial und schl�gt zwei Fligen mit einer Klappe:
 \begin{enumerate}
	 \item wir w�rden verstehen, was das \textcolor{blue}{Spektrum} von \(f\) ist,
	 \item wir h�tten eine der Fourier-Synthese �hnlich Darstellung von \(f\) gefunden.
 \end{enumerate}
 Leider handelt sich bei dem Grenz�bergang \(T\to \infty\) um eine mathematisch heikle Angelegenheit: die Variable \(T\) steckt ja �berall drin: im Integranden und in den Integrationsgrenzen. Der Grenz�bergang klappt f�r Funktionen, die die Voraussetzungen (a) und (b) erf�llen. Um die Details zu verstehen braucht man mehr mathematische Kenntnisse als wir sie zur Verf�gung haben. Versuchen wir trotzdem ein grunds�tzliches Verst�ndnis daf�r zu entwickeln, was f�r \(T\to \infty\) passiert.

 Offensichtlich wird f�r \(T\to\infty\) der Abstand zwischen zwei Kreisfrequenzen beliebig klein! Wir l�sen die Bindung an eine diskrete Variable \(\omega_k\) auf, indem wir auf den Laufindex \(k\) verzichten und das geht so:
 \begin{itemize}
	 \item Die Kreisfrequenz \(\omega = \omega(k)\) wird als Laufvarible in der Summe benutzt und auf den ganzzahligen Laufindex wird komplett verzichtet,
	 \item Anstatt \(k\) f�hren die Weite \(\Delta\omega\) zwischen zwei Frequenzen explizit ein \[ \Delta\omega = \omega_{k+1} - \omega_k = \frac{2\pi}{T} \; \Rightarrow \; \frac{1}{T} = \frac{\Delta \omega}{ 2\pi} \,.\]
 \end{itemize}
 Damit
\[\begin{array}{l rcl}
\forall t \in \left(-\frac{T}{2}, \frac{T}{2}\right) : & f(t) &=& \displaystyle\sum\limits_{k=-\infty}^\infty \left( \dfrac{1}{T} \,  \displaystyle \int\limits_{-\frac{T}{2}}^{\frac{T}{2}} f(t) \cdot  e^{-i \, \omega_k \, t}\, \mathrm{d}t \right) \cdot  e^{i \, \omega_k \, t} \\
& &=& \displaystyle\sum\limits_{\omega=-\infty}^\infty \left( \dfrac{\Delta\omega}{2\pi} \,  \displaystyle \int\limits_{-\frac{T}{2}}^{\frac{T}{2}} f(t) \cdot  e^{-i \, \omega \, t}\, \mathrm{d}t \right) \cdot  e^{i \, \omega \, t} \\
& &=& \displaystyle\sum\limits_{\omega=-\infty}^\infty \dfrac{1}{\sqrt{2\pi}} \left(\dfrac{1}{\sqrt{2\pi}} \,  \displaystyle \int\limits_{-\frac{T}{2}}^{\frac{T}{2}} f(t) \cdot  e^{-i \, \omega \, t}\, \mathrm{d}t \right) \cdot  e^{i \, \omega \, t} \cdot \Delta\omega\,,
\end{array} \]
das Summensymbol bedeutet, dass $\omega$ schrittweise um $\Delta \omega$ nach \(+\infty\) inkrementiert bzw nach \(-\infty\) dekrementiert wird. Schickt man \(T\to \infty\), so wird aus der Summe ein Integral �ber die stetige Variable \(\omega\). Es gilt
\[
\lim\limits_{T\to\infty} \,\sum\limits_{\omega=-\infty}^\infty \left( \dfrac{1}{2\pi} \int\limits_{-\frac{T}{2}}^{\frac{T}{2}} e^{-i\omega t} f(t) \, \mathrm{d}t\right) \cdot e^{i\omega t} \cdot \Delta \omega = \int\limits_{-\infty}^\infty \left( \dfrac{1}{2\pi} \int\limits_{-\infty}^{\infty} e^{-i\omega t} f(t) \, \mathrm{d}t\right)  e^{i\omega t} \,\mathrm{d}\omega
\] und \[ \begin{array}{rcl} \forall t \in \mathbb R: \quad f(t) &=& \displaystyle \int\limits_{-\infty}^\infty \left( \dfrac{1}{2\pi}\displaystyle \int\limits_{-\infty}^{\infty} e^{-i\omega t} f(t) \, \mathrm{d}t\right)  e^{i\omega t} \,\mathrm{d}\omega \\ &=& \dfrac{1}{\sqrt{2\pi}} \, \displaystyle \int\limits_{-\infty}^\infty \left( \textcolor{blue}{ \dfrac{1}{\sqrt{ 2\pi}} \, \displaystyle \int\limits_{-\infty}^{\infty} e^{-i\omega t} f(t) \, \mathrm{d}t } \right)  e^{i\omega t} \,\mathrm{d}\omega =  \dfrac{1}{\sqrt{2\pi}} \, \displaystyle \int\limits_{-\infty}^\infty \textcolor{blue}{ \hat f(\omega) } \cdot e^{i\omega t} \,\mathrm{d}\omega \,. \end{array} \]
Der \textcolor{blue}{ blau eingef�rbte Term} ist die Fourier-Transformierte \( \hat f: \omega \to \hat f(\omega)\). Die Fourier-Transformierte stellt das stetige Spektrum zu \(f\) dar.

Die (kontinuierliche) Fourier-Transformierte $\hat f$ einer Funktion $f$ (zu Einschr�nkungen siehe den vorigen Abschnitt) wird definiert durch:
\begin{defbox}{Fourier-Transformation und inverse Fourier-Transformation}{ }
	\begin{itemize}
		\item Die Abbildung \(\mathcal F: f\to \hat f\) nennt man die Fouiertransformation [Fourier transform]. Die Fouriertransformation ist ein Basiswechsel: vom Zeitbereich in den Frequenzbereich.
		\item Die Fourier-Transformierte \(\hat f\) zu einer Funktion $f$ berechnet sich wie folgt:
	\[
		\hat f(\omega) = \frac{1}{\sqrt{2\pi}} \int\limits_{-\infty}^{\infty} f(t)\cdot e^{-i\, \omega\, t}\, \mathrm{d}t\,.
	\]
\item Die Abbildung \(\mathcal F^{-1}: \hat f\to f\) nennt man die inverse Fouiertransformation [inverse Fourier transform]. Die inverse Fouriertransformation ist ein Basiswechsel: vom Frequenzbereich in den Zeitbereich.
\item Die Funktion \(f\) kann man bei gegebener Fourier-Transformierten \(\hat f\) wie folgt berechnen:
	\[
		f(t) = \frac{1}{\sqrt{2\pi}} \int\limits_{-\infty}^{\infty} \hat f(\omega)\cdot  e^{i\,\omega\, t} \, \mathrm{d}\omega\,.
	\]
	\end{itemize}
\end{defbox}

\noindent Bemerkungen:
\begin{itemize}
	\item �hnlich wie bei der Berechnung der Fourierkoeffizienten, ist der Fall \(\omega = 0 \) separat zu betrachten: \[ \begin{cases} \omega=0:\quad \hat f(0) &= \dfrac{1}{\sqrt{2\pi}} \displaystyle{ \int\limits_{-\infty}^\infty } f(t) \, \mathrm{d} t\,, \\ \omega\not=0:\quad \hat f(\omega) &= \dfrac{1}{\sqrt{2\pi}} \displaystyle{ \int\limits_{-\infty}^\infty } f(t) \cdot e^{-i\,\omega\, t}\,\mathrm{d}t\,. \end{cases}\] Mit einem Grenzprozess l��t sich verifizieren, dass gilt \[ \lim\limits_{\omega\to 0} \hat f(\omega) = \hat f(0)\,.\]
	\item Wie die Fourier-Analyse eines periodischen Signals ist die Fourier-Transformation eine Umwandlung des Signal von Zeit- in Frequenzdarstellung (Basiswechsel).
	\item Die inverse Fourier-Transformation macht das Umgekehrte. Anders als bei periodischen Signalen ist die Synthese vollkommen �quivalent zur Analyse! (Bei periodischen Signalen ist die Fourier-Synthese das Rekonstruieren des Signals im Zeitbereich mittels Fourier-Reihe.)
	\item In der Literatur sieht man die Fourier-Transformierte auch mit anderen Vorfaktoren. %ist der Faktor $\frac12 \pi$ nicht immer per Quadratwurzel zu gleichen Teilen in die Fourier-Transformation und die inverse Fourier-Transformation eingebaut. Es k�nnen also andere Vorfaktoren auftauchen.
	\item Und statt $\hat f$ mit Hut sieht man auch $\tilde f$ mit Tilde.
	\item Was ein Echtzeit-Spectrum-Analyzer zeigt,  ist nicht die  Fourier-Transformation, sondern immer nur die Fourier-Reihe eines kurzen, weich abgeschnittenen Ausschnitts des Signals, praktisch immer per FFT berechnet.
	\item Die Fourier-Transformation ist eine lineare Abbildung:\marginpar{\small Aufgabe}
		 \[
			 \begin{cases} \mathcal F(f+g) = \mathcal F(f) + \mathcal F(g) \,,\\[2ex]
				 \mathcal F(\lambda\cdot f) = \lambda\cdot \mathcal F(f) \,.
		 \end{cases}
		 \]
	 \item Die Fourier-Transformation ver�ndert die L�nge nicht, sie verh�lt sich wie eine Drehung. Es gilt der Satz von Plancherel (entsprechendes gilt f�r die Fourier-Reihe und hei�t der Satz von Parseval):\marginpar{\small Aufgabe}
		\[
			\|f\|^2 = \|\mathcal F(f)\|^2 \Leftrightarrow \langle f, f \rangle = \langle \hat f, \hat f\rangle  \Rightarrow \int_{-\infty}^{\infty} |f(t)|^2 \,\mathrm{d}t = \int_{-\infty}^{\infty} |\hat f(\omega)|^2 \,\mathrm{d}\omega \,.
		\]
Praktisch bedeutet dies, dass man die physikalische Arbeit sowohl in der Zeit- als auch Frequenzdarstellung berechnen kann.
\end{itemize}

\begin{exabox}{Fourier-Transformation}{}
	Funktion:
	\[ f:\mathbb R \to \mathbb R, t\mapsto e^{-|t|} \]
	Fourier-Transformation:
	\begin{eqnarray*}
		\omega=0: \quad \hat f(0) &=& \frac{1}{\sqrt{2\pi}}\int\limits_{-\infty}^\infty f(t) \, \mathrm{d} t =  \frac{1}{\sqrt{2\pi}}\int\limits_{-\infty}^0 e^t \, \mathrm{d} t +  \frac{1}{\sqrt{2\pi}}\int\limits_{0}^\infty e^{-t} \, \mathrm{d} t \\
	&=&  \frac{1}{\sqrt{2\pi}}\Big[ e^t \Big]_{-\infty}^0 +  \frac{1}{\sqrt{2\pi}}\Big[-e^{-t} \Big]_{0}^\infty  =  \frac{1}{\sqrt{2\pi}}\left( 1 - 0 \right) +  \frac{1}{\sqrt{2\pi}}\left( 0 + 1 \right)   = \frac{2}{\sqrt{2\pi}}\,.
\end{eqnarray*}
	\begin{eqnarray*}
	\omega\not=0:\quad \hat f(\omega) &=& \frac{1}{\sqrt{2\pi}}\int\limits_{-\infty}^\infty f(t) \cdot e^{-i\omega t} \, \mathrm{d} t =  \frac{1}{\sqrt{2\pi}}\int\limits_{-\infty}^0 e^t \cdot e^{-i\omega t} \, \mathrm{d} t +  \frac{1}{\sqrt{2\pi}}\int\limits_{0}^\infty e^{-t} \cdot e^{-i\omega t} \, \mathrm{d} t \\
	&=&  \frac{1}{\sqrt{2\pi}}\Big[\frac{e^{(1-i\omega)t}}{1-i\omega} \Big]_{-\infty}^0 +  \frac{1}{\sqrt{2\pi}}\Big[\frac{e^{-(1+i\omega)t}}{-(1+i\omega)} \Big]_{0}^\infty \\
	&=&  \frac{1}{\sqrt{2\pi}} \frac{1}{1-i\omega}\left( 1 - \underbrace{\lim\limits_{t\to-\infty}e^{(1-i\omega)t}}_{\displaystyle{\stackrel{(*)}{=}} 0 } \right)-  \frac{1}{\sqrt{2\pi}} \frac{1}{1+i\omega}\left( \underbrace{\lim\limits_{t\to \infty} e^{-(1+i\omega)t}}_{\displaystyle{\stackrel{(*)}{=} 0}} - 1\right)  \\
		&=&  \frac{1}{\sqrt{2\pi}}\frac{1}{1-i\omega} +  \frac{1}{\sqrt{2\pi}}\frac{1}{1+i\omega}  =  \frac{1}{\sqrt{2\pi}}\frac{2}{1+\omega^2} = \sqrt{\frac{2}{\pi}} \frac{1}{1+\omega^2}\,,
	\end{eqnarray*}
	wobei die Grenzwerte in \( (*)\)  null sind, weil die komplexe Exponentialfunktion beschr�nkt ist und der reelle Anteil gegen null konvergiert: \[ \lim\limits_{t\to\pm\infty} \left| e^{\mp t \mp i \omega t} \right| \leq \lim\limits_{t\to\infty}e^{-|t|} \cdot 1  = 0\,. \]
\end{exabox}


Anders als bei der Fourier-Reihe entwickelt man seltener die Fourier-Transformierte in reeller Form. F�r Signale mit Symmetrieeigenschaften kann das aber sinnvoll sein. Es gilt\marginpar{\small Aufgabe}

\begin{defbox}{Fourier-Transformation f�r Signale mit Symmetrieeigenschaften}{ }
	\begin{itemize}
		\item[(a)] Die Fourier-Transformierte einer Funktion $f$ mit \( f(t) = f(-t)\) lautet:
	\[
		\hat f(\omega) = \frac{2}{\sqrt{2\pi}} \int\limits_{0}^{\infty} f(t) \cdot \cos(\omega \, t) \, \mathrm{d}t\,.
	\]
		\item[(b)] Die Fourier-Transformierte einer Funktion $f$ mit \( f(t) = -f(-t)\) lautet:
	\[
		\hat f(\omega) = \frac{-2\,i}{\sqrt{2\pi}} \int\limits_{0}^{\infty} f(t) \cdot \sin(\omega \, t) \, \mathrm{d}t\,.
	\]
	\end{itemize}
\end{defbox}
